\section{A local potential from a curve and the curvature density}\label{sec:potential}

Let $X$ be a smooth projective surface, $H$ an ample line bundle and
$\omega=c_1(H,h_H)$ a K\"ahler form. We use the complex orientation and the
conventions
\begin{equation}\label{eq:conventions}
 \ddc=\frac{i}{\pi}\partial\bar\partial,\qquad
 c_1(H,h_H)=\frac{i}{2\pi}F_{h_H},\qquad
 \Vol(X)=\frac12\int_X\omega^2.
\end{equation}
For a unitary connection $A$ on a rank-$r$ bundle $E$, write
$F_A^0=F_A-r^{-1}\tr(F_A)\,\mathrm{id}$.
The determinant connection is fixed. Tangent vectors to the smooth moduli
space of irreducible Hermitian--Yang--Mills connections are represented by
trace-free skew-Hermitian one-forms in Coulomb gauge, and
\begin{equation}\label{eq:l2metric}
 \gL(a,b)=-\int_X\tr(a\wedge *b),\qquad
 \Om(a,b)=\gL(Ja,b).
\end{equation}
Here $(Ja)^{0,1}=ia^{0,1}$. On an obstructed moduli space all differential
identities below are statements on its smooth stable locus, or after
pullback to a smooth holomorphic family of stable bundles.
The trace is the trace in the defining representation. For a Hermitian--Yang--Mills connection,
\begin{equation}\label{eq:charge}
 \int_X|F_A^0|^2\,dV_X=8\pi^2\kappa,\qquad
 \kappa=c_2(E)-\frac{r-1}{2r}c_1(E)^2.
\end{equation}
The trace-free curvature is primitive of type $(1,1)$ and hence
anti-self-dual. One point bubble of charge one has curvature mass $8\pi^2$.
These conventions fix the constants both in the potential and in the
metric induced on a boundary stratum.

\subsection{The fibre integral}
Let $p:S\times X\to S$ and $q:S\times X\to X$ be the projections, where
$S$ is a complex manifold. Suppose that $\mathbb E$ is a holomorphic family
of stable bundles with a smooth Hermitian metric that is Hermitian--Einstein
on each fibre. Denote its classifying map by $f:S\to\M$ and set
\[
 \Theta=\ch_2(\End\mathbb E,h).
\]
All fibre integrals below use the complex orientation.

\begin{lemma}\label{lem:fibre-integral}
The pullback of the $L^2$ K\"ahler form is
\begin{equation}\label{eq:fibre-integral}
 f^*\Om=-\frac{2\pi^2}{r}\,p_*(\Theta\wedge q^*\omega).
\end{equation}
\end{lemma}
\begin{proof}
Let $F_0$ be the trace-free curvature of the Chern connection on the family.
The identity
$\tr_{\End}(\operatorname{ad}a\operatorname{ad}b)=2r\tr(a_0b_0)$ gives
\[
 \Theta=-\frac{r}{4\pi^2}\tr(F_0\wedge F_0).
\]
For $v\in T_sS$ put $a_v=(\iota_vF_0)|_{X_s}$.
The mixed curvature of a fibrewise Hermitian--Einstein family represents
the harmonic Kodaira--Spencer class; this follows by differentiating the
Hermitian--Einstein equation and using the Bianchi and K\"ahler identities
\cite[Proposition~1]{ST92}. Thus $a_v$ is the Coulomb representative of
$df(v)$ and $a_{Iv}=Ja_v$.
For real vectors $v,w$, contraction gives
\[
 \iota_w\iota_v\tr(F_0\wedge F_0)
 =2\tr\big(F_0(v,w)F_0\big)
  -2\tr\big(\iota_vF_0\wedge\iota_wF_0\big).
\]
The first term vanishes after restriction to $X_s$ and multiplication by
$\omega$, because $F_0|_{X_s}\wedge\omega=0$. Consequently
\[
 p_*(\Theta\wedge q^*\omega)(v,w)
 =\frac{r}{2\pi^2}\int_X\tr(a_v\wedge a_w)\wedge\omega.
\]
Pointwise, $\tr(a\wedge Ja)\wedge\omega=-|a|^2dV_X$.
Polarization identifies the integral on the right with $-\Om(df(v),df(w))$,
which proves \eqref{eq:fibre-integral}.
\end{proof}

The fibre-integral description is classical; compare
\cite[Proposition~4.3]{BS16}. The calculation above records its coefficient
for the real metric \eqref{eq:l2metric}.

\subsection{Transgression to an auxiliary curve}
Choose a smooth curve $C\in|H^k|$ and a section $s_C$ defining it. Put
\begin{equation}\label{eq:green}
 g_C=\frac1k\log|s_C|_{h_H^k}.
\end{equation}
The logarithm is of the norm, not its square. Poincar\'e--Lelong gives
\begin{equation}\label{eq:pl}
 \ddc g_C=\frac1k[C]-\omega.
\end{equation}
In particular, $g_C\in L^1(X)$ and $\ddc(-g_C)=\omega$ on $X\setminus C$.

\begin{lemma}\label{lem:transgression}
If $\Theta$ is a smooth closed $(2,2)$-form on $S\times X$, then
\begin{equation}\label{eq:transgression}
 p_*(\Theta\wedge q^*\ddc g_C)
 =\ddc\left(\int_Xg_C\Theta|_{X_s}\right).
\end{equation}
The integral on the right is a smooth function of $s$.
\end{lemma}
\begin{proof}
Since $\Theta$ has pure type and is closed, $d^c\Theta=0$.
Multiplication of a distribution by a smooth form therefore gives
$\ddc(q^*g_C\Theta)=q^*\ddc g_C\wedge\Theta$.
Fibre integration commutes with $\partial$ and $\bar\partial$.
The asserted smoothness follows by differentiating the smooth coefficients
of $\Theta$ under the integral; on compact subsets of $S$, their derivatives
are bounded and the fixed function $g_C$ is integrable.
\end{proof}

On a Hermitian--Yang--Mills fibre,
\[
 \Theta|_{X_s}=-\frac{r}{4\pi^2}|F_{A_s}^0|^2dV_X.
\]
Substitution of \eqref{eq:pl} in \eqref{eq:fibre-integral} consequently yields
\begin{equation}\label{eq:curve-transgression}
 f^*\Om=-\frac{2\pi^2}{rk}(p_C)_*(\Theta|_{S\times C})
       +\ddc\left(-\frac12\int_Xg_C|F_{A_s}^0|^2dV_X\right).
\end{equation}
This is the step that replaces a potential on the surface by a determinant
on a curve and a scalar integral on the surface.

\subsection{The determinant term}
Let $G$ be a Hermitian holomorphic bundle of rank $\rho$ on $C$, and suppose
\begin{equation}\label{eq:acyclic}
 H^j(C,\End E_s|_C\otimes G)=0\quad(j=0,1)
\end{equation}
for $s\in S$. Fix a metric on $C$. Let $\Delta_s=\bar\partial^*\bar\partial$
on sections of $\End E_s|_C\otimes G$ with the induced metric.
It is positive and invertible. Its zeta determinant is
$\det_\zeta\Delta_s=\exp(-\zeta'_{\Delta_s}(0))$.

\begin{lemma}\label{lem:quillen-curve}
Under \eqref{eq:acyclic},
\begin{equation}\label{eq:quillen-curve}
 (p_C)_*(\Theta|_{S\times C})
   =\frac1{2\rho}\ddc\log\det_\zeta\Delta_s.
\end{equation}
\end{lemma}
\begin{proof}
Apply the curvature formula of Bismut--Gillet--Soul\'e to
$\xi=\End\mathbb E|_{S\times C}\otimes q^*G$
\cite[Theorem~1.6, Remark~1.10 and equation~(1.154)]{BGS88}.
In the acyclic case the squared Quillen norm of the canonical determinant
section is $\det_\zeta\Delta_s$, with
\[
 \frac{i}{2\pi}\partial\bar\partial\log\det_\zeta\Delta_s
 =\left[(p_C)_*\big(\ch(\xi)\operatorname{Td}(C)\big)\right]^{1,1}.
\]
The Chern form $\ch_1(\End\mathbb E,h)$ vanishes identically.
All forms involving $G$ and $\operatorname{Td}(C)$ are pulled back from
$C$. Thus the only term contributing a $(1,1)$-form on $S$ is
$\rho(p_C)_*\Theta$. The convention \eqref{eq:conventions} gives the factor
$2$ in \eqref{eq:quillen-curve}.
\end{proof}

\begin{proposition}\label{prop:potential}
On the locus \eqref{eq:acyclic}, the function
\begin{equation}\label{eq:potential}
 u(A_s)=-\frac{\pi^2}{rk\rho}\log\det_\zeta\Delta_s
        -\frac12\int_Xg_C|F_{A_s}^0|^2dV_X
\end{equation}
is smooth, gauge invariant, and satisfies $\ddc u=f^*\Om$.
\end{proposition}
\begin{proof}
Combine \eqref{eq:curve-transgression} and \eqref{eq:quillen-curve}.
Both the determinant and the curvature integral are invariant under unitary
gauge transformations. Tensoring a universal family by a line bundle from
$S$ leaves its endomorphism bundle unchanged. Thus the same function is
defined on the corresponding locus of gauge classes.
\end{proof}

The metric on $E_s|_C$ varies with $s$. In particular, positivity of a
Quillen line for a fixed curve metric does not separately imply
plurisubharmonicity of the first term of \eqref{eq:potential}.
Positivity is asserted for their sum, whose complex Hessian is \eqref{eq:l2metric}.

For a K\"ahler form with $[\omega]=t\,c_1(H)$, $t>0$, choose
$c_1(H,h_H)=\omega/t$ and replace \eqref{eq:green} by
$g_C=(t/k)\log|s_C|$.
The determinant coefficient in \eqref{eq:potential} becomes
$-t\pi^2/(rk\rho)$; the curvature coefficient remains $-1/2$.
Thus the formulas apply on any positive ray of a polarization with the
expected scaling. We use $t=1$ in the statements below.
